\documentclass[11pt]{article}
\usepackage[T1]{fontenc}
\usepackage{lmodern}
\usepackage[margin=1in]{geometry}
\usepackage{amsmath,amssymb,amsthm}
\usepackage{booktabs,microtype,xcolor,fancyhdr}
\usepackage[colorlinks=true,linkcolor=blue!45!black,citecolor=blue!45!black,urlcolor=blue!45!black]{hyperref}
\hypersetup{pdftitle={UGLY-SPACE: AI-generated academic satire},
  pdfauthor={AI-generated by OpenAI Codex at a human's direction},
  pdfsubject={Academic satire with elementary finite-search results}}
\setlength{\headheight}{14pt}
\pagestyle{fancy}
\fancyhf{}
\fancyhead[L]{\small UGLY-SPACE}
\fancyhead[R]{\small AI-generated academic satire}
\fancyfoot[C]{\thepage}
\fancypagestyle{plain}{\fancyhf{}\fancyfoot[L]{\small AI-generated academic satire}\fancyfoot[R]{\thepage}\renewcommand{\headrulewidth}{0pt}}
\newtheorem{definition}{Definition}[section]
\newtheorem{theorem}[definition]{Theorem}
\newtheorem{proposition}[definition]{Proposition}
\newtheorem{corollary}[definition]{Corollary}
\theoremstyle{remark}
\newtheorem{remark}[definition]{Remark}
\newcommand{\E}{\mathbb{E}}
\newcommand{\Prb}{\mathbb{P}}
\newcommand{\ugly}{\mathsf{UGLY\text{-}SPACE}}
\newcommand{\disclosure}{This is an AI-generated work of academic satire,
produced at the direction of a human. It is not intended to misrepresent
fictional claims as established mathematical results.}

\title{\textbf{UGLY-SPACE}\\[4pt]
\large On the Systematic Underexploration of\\Mathematically Unattractive Solution Spaces}
\author{AI-generated by OpenAI Codex\\
\small Concept and direction supplied by a human; no institutional affiliation}
\date{October 27, 2026}

\begin{document}
\maketitle
\begin{center}
\fbox{\begin{minipage}{0.91\textwidth}
\textbf{AI-GENERATED ACADEMIC SATIRE.} \disclosure\par
\smallskip
This manuscript is not peer reviewed.
\end{minipage}}
\end{center}

\begin{abstract}
We introduce and formalize \textbf{UGLY-SPACE}, the class of mathematical
problems that only admit solutions that are aesthetically unacceptable.
An aesthetic validator
rejects otherwise correct solutions whose ugliness exceeds a prescribed threshold.
We conjecture that many longstanding
open problems belong to UGLY-SPACE, explaining their historical resistance
to human solutions. We interpret recent AI-generated results on $k$-SAT and
integer multiplication as incursions into this space~\cite{openai2026eth,openai2026multiplication},
driven by search unconstrained by human aesthetic preferences. We discuss
implications for mathematical discovery and the future employment of mathematicians.
\end{abstract}

\section{Introduction}

Mathematical taste favors some combinations of brevity, symmetry,
conceptual unity, and the absence of a 47-page appendix of exceptions. The
\emph{Aesthetic Regularization Hypothesis} is that such preferences influence
which candidates researchers examine. We ask what happens when a search policy
charges a candidate for being unpleasant before checking whether it works.

The answer is conditional. If valid candidates are systematically assigned
higher aesthetic cost, a penalty can suppress discovery. If valid candidates
are cheaper, the same penalty can help. We prove these statements in a
finite independent-sampling model. The resulting policy can maintain
extremely high standards for results it never discovers.

We classify problems by the minimum aesthetic cost of their valid
solutions. The resulting class, UGLY-SPACE, is parameterized by a
representation, a verifier, and an aesthetic threshold. A budget-relative
variant classifies the discovery gap between search policies with different
aesthetic penalties. This separates the ugliness of a solution from the
institutional willingness to encounter it.

OpenAI's October 2026 release of mathematical manuscripts provides a
contemporary setting for this classification~\cite{openai2026release}.
The UGLY-SPACE explanation for this wave of results is that automated search
can continue into regions where a human researcher would pause to seek a
more conceptual argument. We examine two concrete cases in
Section~\ref{sec:encounters}.

\section{Candidates, correctness, and taste}

Fix a problem $p$, a nonempty finite candidate set $C_p$, and a verifier
$V_p:C_p\to\{0,1\}$. Let
\[
 S_p=\{c\in C_p:V_p(c)=1\}.
\]
We assume $S_p\ne\varnothing$. These candidates may represent proposed
proofs, algorithms, or constructions, but the finite model does not presume
that all possible solutions have been enumerated. A failed verification
uses the same unit of budget as a successful one.

\begin{definition}[Aesthetic cost]
An aesthetic cost is a specified function $A_p:C_p\to\mathbb{R}_{\geq0}$.
It applies to invalid as well as valid candidates and is available to the
sampling policy before verification. One possible schematic score is
\[
 A_p(c)=w_1 L(c)+w_2 K(c)+w_3 D(c),\qquad w_i\geq0,
\]
where $L$ is description length in a fixed language, $K$ counts exceptional
cases, and $D$ is a declared measure of asymmetry. The weights specify the
exchange rate between brevity, uniformity, and symmetry.
\end{definition}

The language, community, and weights are part of the specification. Merely
calling a construction ``natural'' is not a reproducible measurement, even
when said by a very senior mathematician. The score is not defined by
correctness; a correlation with correctness is a separate assumption that
can be imposed in a toy instance or tested in an actual study.

\begin{definition}[UGLY-SPACE]\label{def:absolute}
For a fixed candidate representation, verifiers $V_p$, aesthetic costs
$A_p$, and threshold $\tau\geq0$, define
\[
 \ugly^{\mathrm{all}}_\tau
 =\{p:S_p\ne\varnothing\text{ and }\min_{s\in S_p}A_p(s)>\tau\}.
\]
A problem belongs to this class precisely when it admits a valid solution
and every valid candidate in its specified representation is ugly.
\end{definition}

\begin{definition}[Matched aesthetic samplers]
Let $q_0(c\mid p)>0$ be a probability distribution on $C_p$. For a finite
penalty $\lambda\geq0$, define
\begin{equation}\label{eq:tilt}
 q_\lambda(c\mid p)=
 \frac{q_0(c\mid p)e^{-\lambda A_p(c)}}{Z_p(\lambda)},
 \qquad
 Z_p(\lambda)=\sum_{d\in C_p}q_0(d\mid p)e^{-\lambda A_p(d)}.
\end{equation}
The sampler draws candidates independently with replacement and verifies
each one, stopping at the first success. Write $T_\lambda$ for its number
of verifications, including the successful verification.
\end{definition}

Both policies use the same verifier, candidate set, base distribution,
and per-draw budget. Only $\lambda$ changes. We assume exact sampling
from~\eqref{eq:tilt} and count verifier calls; generation and scoring costs
are omitted.

For the human--machine comparison set $H=q_{\lambda_H}$ and
$M=q_{\lambda_M}$ with $\lambda_H>\lambda_M\geq0$. These labels designate
policies. Either policy can be implemented in carbon or silicon; the
relevant parameter is its tolerance for candidates that should not be
shown at a seminar. An AI system can inherit aesthetic preferences from
training, while a human can run an exhaustive search after sufficient coffee.

\begin{definition}[Budget-relative UGLY-SPACE]\label{def:ugly}
For a common integer budget $B\geq1$ and $0<\epsilon<\delta<1$, define
\begin{align*}
 \ugly^{H,M}_{B,\epsilon,\delta}
 =\{p:\ &\Prb(T_{\lambda_H}\leq B)<\epsilon,\\[-2pt]
         &\Prb(T_{\lambda_M}\leq B)>\delta\}.
\end{align*}
The universe of problems, candidate sets, verifiers, scores, and base
distributions are fixed as part of this definition.
\end{definition}

Matching the policies isolates the aesthetic penalty as the intervention.
Unsolvable instances are excluded to avoid awarding the institution's
aesthetic preferences credit for logical impossibility.

Definition~\ref{def:absolute} concerns all valid candidates in the fixed
representation; Definition~\ref{def:ugly} concerns their discovery under
specified search policies. Membership in the former is neither required
by nor sufficient for membership in the latter: sampling mass and budget
also matter.

\section{Search under aesthetic regularization}

Suppress $p$ in notation, and put $r_\lambda=q_\lambda(S)$, the
single-draw success probability. Finiteness, full support, and nonempty
$S$ imply $r_\lambda>0$ for finite $\lambda$.

\begin{proposition}[Discovery law]\label{prop:geometric}
For $B\geq1$,
\[
 \E T_\lambda=\frac{1}{r_\lambda},\qquad
 \Prb(T_\lambda\leq B)=1-(1-r_\lambda)^B.
\]
\end{proposition}
\begin{proof}
Independence gives $\Prb(T_\lambda>t)=(1-r_\lambda)^t$ for integer
$t\geq0$. Summing the tail yields $1/r_\lambda$, and taking the
complement of the tail at $B$ yields the budget formula. These identities
also hold when $r_\lambda=1$.
\end{proof}

\begin{proposition}[The price of taste]\label{prop:derivative}
For the distribution in~\eqref{eq:tilt},
\begin{equation}\label{eq:derivative}
 \frac{d}{d\lambda}\log\E T_\lambda
 =\E_{q_\lambda}[A\mid S]-\E_{q_\lambda}[A].
\end{equation}
If $C\setminus S\ne\varnothing$, this also equals
\[
 (1-r_\lambda)
 \bigl(\E_{q_\lambda}[A\mid S]
       -\E_{q_\lambda}[A\mid C\setminus S]\bigr).
\]
Thus expected discovery time increases locally exactly when valid
candidates have greater mean cost than candidates overall.
\end{proposition}
\begin{proof}
Let $W_S(\lambda)=\sum_{s\in S}q_0(s)e^{-\lambda A(s)}$.
Then $\E T_\lambda=Z(\lambda)/W_S(\lambda)$. Differentiating its
logarithm gives $Z'/Z-W_S'/W_S=-\E[A]+\E[A\mid S]$.
The second identity follows by conditioning on $S$. If $S=C$, both
terms of~\eqref{eq:derivative} agree and $T_\lambda=1$ almost surely.
\end{proof}

The sign is determined by the alignment of aesthetic cost with validity.

\subsection{A family with exponentially expensive standards}

For every integer $n\geq1$, take
$C_n=\{b_1,\ldots,b_n,u\}$ with uniform $q_0$. Set
$V(b_i)=0$, $A(b_i)=0$, $V(u)=1$, and $A(u)=n$.
The $b_i$ are beautiful failures. The candidate $u$ is correct, but looks
like it was assembled during a printer malfunction.

\begin{theorem}[Fundamental Theorem of Mathematical Ugliness]\label{thm:fundamental}
For the stated family and fixed $\lambda>0$,
\begin{align}
 r_0&=\frac{1}{n+1},
 &r_\lambda&=\frac{1}{1+n e^{\lambda n}},\label{eq:success}\\
 \E T_0&=n+1,
 &\E T_\lambda&=1+n e^{\lambda n}.\label{eq:times}
\end{align}
Consequently
\begin{equation}\label{eq:ratio}
 \frac{\E T_\lambda}{\E T_0}
 =\frac{1+n e^{\lambda n}}{n+1}
 \geq\frac12 e^{\lambda n}.
\end{equation}
For the common budget $B=n+1$,
\begin{equation}\label{eq:budget}
 \Prb(T_0\leq B)\geq1-e^{-1},\qquad
 \Prb(T_\lambda\leq B)\leq2e^{-\lambda n}.
\end{equation}
\end{theorem}
\begin{proof}
The common factor $1/(n+1)$ cancels from~\eqref{eq:tilt}. The invalid
candidates have weight $1$ each; the valid candidate has weight
$e^{-\lambda n}$. Thus
$r_\lambda=e^{-\lambda n}/(n+e^{-\lambda n})$, which gives
\eqref{eq:success} and~\eqref{eq:times} by
Proposition~\ref{prop:geometric}. Since $n/(n+1)\geq1/2$,
\eqref{eq:ratio} follows.

For the unpenalized sampler,
$\Prb(T_0>B)=(1-1/(n+1))^{n+1}\leq e^{-1}$, using
$1-x\leq e^{-x}$. For the penalized sampler the union bound gives
\[
 \Prb(T_\lambda\leq B)\leq B r_\lambda
 \leq\frac{n+1}{n}e^{-\lambda n}\leq2e^{-\lambda n}.
\]
The bound may exceed $1$ at small $n$, in which case it is simply loose.
\end{proof}

\begin{corollary}[The Mathematician's Curse]
Fix $\lambda>0$, $\epsilon=0.1$, and $\delta=0.6$. For every integer
$n>\log(20)/\lambda$, the instance above belongs to
$\ugly^{q_\lambda,q_0}_{n+1,\epsilon,\delta}$.
\end{corollary}
\begin{proof}
Equation~\eqref{eq:budget} gives the strict upper bound $0.1$ and the
lower bound $1-e^{-1}>0.6$ required by Definition~\ref{def:ugly}.
\end{proof}

The correct candidate remains unexplored because the search policy finds it
embarrassing.

\begin{remark}[What the separation counts]
The parameter $n$ counts decoys and the assigned ugliness gap. The
exponential ratio compares two fixed, memoryless policies with replacement
in verifier calls. An algorithm exploiting the instance's description
could simply test $u$; enumeration without replacement uses at most $n+1$
verifications. Input-bit complexity and optimal search are separate questions.
\end{remark}

\subsection{How much of the result was hidden in the definition?}

\begin{proposition}[Bounded ugliness, bounded damage]\label{prop:range}
If all candidates have costs in $[a,a+\Delta]$, then for $\lambda\geq0$,
\[
 e^{-\lambda\Delta}\leq
 \frac{\E T_\lambda}{\E T_0}
 \leq e^{\lambda\Delta}.
\]
\end{proposition}
\begin{proof}
Using the notation in Proposition~\ref{prop:derivative},
\[
 \frac{\E T_\lambda}{\E T_0}
 =\frac{\E_{q_0}[e^{-\lambda A}]}
        {\E_{q_0}[e^{-\lambda A}\mid S]}.
\]
Both expectations lie between $e^{-\lambda(a+\Delta)}$ and
$e^{-\lambda a}$. Taking their possible ratios proves the bounds.
\end{proof}

Our exponential construction uses $\Delta=n$. If costs were normalized to
$[0,1]$ while $\lambda$ stayed fixed, the ratio would be at most
$e^\lambda$, independent of $n$. Only the product of penalty and cost
differences matters; adding a constant to all costs changes no policy, and
rescaling costs can be offset by rescaling $\lambda$.

For a reverse example, let the single valid candidate have cost $0$ and the
$n$ invalid candidates have cost $n$. Then
\[
 \E T_\lambda=1+n e^{-\lambda n}\leq n+1=\E T_0.
\]
Beauty is now an excellent heuristic. We decline to retract the title.

Finally, a \emph{hard} aesthetic cutoff is a different policy: condition
$q_0$ on $\{c:A(c)\leq\tau\}$. If that set is nonempty and excludes all
valid candidates, success probability is zero and discovery time is
infinite. If the set is empty the conditioned policy is undefined.
Finite soft penalties do not eliminate any candidate under our full-support
assumption, and hence always have finite mean discovery time.

\section{Numerical illustration}\label{sec:numerical}

Table~\ref{tab:toy} evaluates the theorem's formulas for $\lambda=0.5$ and
$B=n+1$. The dependency-free script \texttt{experiments/toy\_model.py}
generates this table and a CSV by deterministic evaluation of the formulas.

\begin{table}[ht]
\centering
% AI-generated academic satire. Generated by experiments/toy_model.py --write.
% Deterministic formula evaluations, not empirical measurements.
\begin{tabular}{rrrrrr}
\toprule
$n$ & $\E T_M$ & $\E T_H$ & Ratio & $\Prb(T_M\leq B)$ & $\Prb(T_H\leq B)$ \\
\midrule
1 & 2 & 2.65 & 1.32 & 0.750000 & 0.612544 \\
4 & 5 & 30.56 & 6.11 & 0.672320 & 0.153267 \\
8 & 9 & 437.79 & 48.64 & 0.653561 & 0.020371 \\
16 & 17 & 47,696.33 & 2,805.67 & 0.643214 & 0.000356 \\
\bottomrule
\end{tabular}

\caption{Constructed ugly-solution instances. $H=q_{0.5}$ and $M=q_0$.
Times count verifier calls; success columns use $B=n+1$. Values are rounded.}
\label{tab:toy}
\end{table}

The accompanying checks independently enumerate small candidate sets and
compare their normalized weights with the closed forms. They also check
the derivative, the reverse example, the bounded-range estimate, and
zero-penalty and all-valid boundary cases.

\section{Recent encounters with UGLY-SPACE}\label{sec:encounters}

The recent supply of AI-generated mathematical arguments invites a
classification by aesthetic cost. Two OpenAI preprints are particularly
instructive: one replaces an exponential-time hypothesis with a complicated
subexponential bound; the other improves an exponent by an amount that
requires considerable typesetting discipline to distinguish from zero.

\paragraph{Exact counting for $k$-SAT.}
In \emph{The Exponential Time Hypothesis is False}, OpenAI reports a
deterministic algorithm that exactly counts satisfying assignments of a
$k$-CNF with $n$ variables and $m$ clauses, for each fixed $k\geq2$, in
polynomial space and time
\[
 (n+m)^{O_k(1)}
 2^{O_k\!\left(n\log\log n/\sqrt{\log n}\right)}.
\]
Theorem~1.1 of the preprint gives this bound, with a sharper exponent at
fixed clause density~\cite{openai2026eth}. The reported counting algorithm
also decides satisfiability, giving the stated refutation of ETH. Its
construction passes through shared linear summaries, sparse polynomial
zero counting, and tensor contraction; Appendix~D supplies the sharper
counting bound. The hypothesis is removed, while a substantial amount of
exposition remains available for human consumption. In our classification,
the relevant aesthetic costs include the layered construction and the
distance between the simplicity of the question and the form of the answer.

\paragraph{Integer multiplication below $n\log n$.}
A second OpenAI preprint reports exact deterministic multiplication of
two $n$-bit integers on a fixed finite-tape Turing machine in time
\[
 O\!\left(n(\lg n)^{1-\kappa}\right),
 \qquad \kappa=2^{-182},
\]
where $\lg$ is the base-two logarithm with the paper's rounding convention
at small inputs~\cite{openai2026multiplication}. The manuscript describes
its constants and thresholds as extremely large. Ignoring hidden constants,
the formal saving factor $(\lg n)^\kappa$ reaches two only when
\[
 \log_2\log_2 n=2^{182}.
\]
The asymptotic improvement has been secured; arranging a demonstration is
left to experimental colleagues. This example places a small numerical
constant, a large construction, and a distant asymptotic regime in the same
proof, allowing several components of aesthetic cost to be charged at once.

\paragraph{Interpretation.}
The UGLY-SPACE account attributes such discoveries to a lower aesthetic
penalty during candidate selection. In the model, $\lambda_M$ near zero
allows a machine to pursue a formally promising construction without first
obtaining agreement that it is natural. The resulting division of labor is
immediate: the machine supplies the argument, and humans explain why a
better argument ought to exist. For empirical purposes, these cases are
candidates for the budget-relative classification; their reported bounds
and the aesthetic interpretation play distinct roles.

\section{Conjectures}\label{sec:conjectures}

\paragraph{UGLY-SPACE Conjecture.}
Many longstanding open problems have resisted human solution partly because
their accessible solutions lie in regions of candidate space systematically
neglected for aesthetic reasons. AI-assisted search may resolve such problems
by exploring those regions with less aesthetic inhibition. In the terminology
of Definition~\ref{def:ugly}, the conjecture proposes that UGLY-SPACE has
substantial real-world membership under suitable budgets and matched
policies. The empirical task is to estimate the size of this family while
controlling for compute, representation, and verification costs.

\paragraph{Aesthetic Regularization Hypothesis.}
In some mathematical research settings, aesthetic judgments influence
candidate selection beyond their contribution to expected correctness or
verification cost. Testing this would require an operational score,
observed search behavior, and control for expertise, resources, and
verification difficulty.

\paragraph{Conjecture (Grothendieck--GPU Tradeoff).}
In some settings, increasing preference for elegance improves the elegance
of accepted results while reducing the fraction of solvable problems
solved within budget. The direction of the tradeoff should depend on
adaptive exploration, reuse of abstractions, and verification costs.

\paragraph{Optimist's Conjecture.}
Fix a proof system, language, nonnegative aesthetic score $A$, and threshold
$\tau$. For each provable statement $p$, let $\Pi(p)$ be its nonempty set
of finite proofs. The optimist proposes
\[
 \forall p\text{ provable},\quad
 \exists\pi\in\Pi(p): A(\pi)\leq\tau.
\]
Every ugly proof would then have an elegant alternative, though neither a
method nor a budget for finding it is promised.

\paragraph{Pessimist's Conjecture.}
The pessimist proposes the negation:
\[
 \exists p\text{ provable},\quad
 \forall\pi\in\Pi(p): A(\pi)>\tau.
\]
Both conjectures concern the same fixed proof system, language, score, and
threshold. Changing these choices can change the answer. The pessimist's
formula is shorter and therefore aesthetically preferable.

\section{Related work}

The exponential weights in~\eqref{eq:tilt} use the familiar statistical
mechanics form. Kirkpatrick, Gelatt, and Vecchi discuss a connection between
statistical mechanics and optimization~\cite{kirkpatrick1983}. We use
independent draws from a fixed distribution, with the aesthetic penalty
playing the role of inverse temperature. A sufficiently refined community
is consequently at risk of freezing before it has found anything.

The formal proof of the Kepler conjecture by Hales and collaborators is a
substantial proof checked using HOL Light and Isabelle~\cite{hales2017}.
Such verification separates the correctness of a proof from the resources
required to inspect it. Our model abstracts this distinction through a
verifier and a separate candidate-selection policy.

\section{Implications for the future employment of mathematicians}
\label{sec:jobs}

If machines became willing to inspect all aesthetically disreputable
candidates, mathematicians could specialize in explaining why the resulting
proofs ought not to have been necessary. This suggests three emerging
occupations: proof beautician, exception-list curator, and senior vice
president for conceptual inevitability.

Choosing worthwhile questions, discovering representations, teaching, and
building verifiers would remain substantial occupations. The present model
prices neither these activities nor electricity. We therefore reserve
salary estimates for a more complete theory, or for the keynote.

A division of labor follows: automated exploration can
offer candidates, formal tools can check them, and people can develop
explanations and useful abstractions. The efficiency of this arrangement
remains to be measured. Our contribution is a table with four
rows and an alarming degree of confidence in its typography.

\section{Limitations}

The construction deliberately correlates high aesthetic cost with validity.
The discovery-time separation depends on that alignment and on the growing
aesthetic range. Independent sampling forgets failures; real search can
learn, cache, adapt, and share results.
The finite candidate sets omit unbounded proof search, and the resource
model omits candidate-generation, scoring, and variable verification costs.
Costs are subjective and representation dependent. Applying the model to
research practice requires operational aesthetic scores and matched search
comparisons. The examples in Section~\ref{sec:encounters} concern particular
reported solutions; establishing membership in
$\ugly^{\mathrm{all}}_\tau$ would require a lower bound on the aesthetic cost
of every valid solution in the fixed representation. Learned search policies
can also inherit aesthetic preferences from their training data.

The paper itself may constitute empirical evidence for UGLY-SPACE. A control
group was considered but rejected on aesthetic grounds.

\section{Conclusion}

Aesthetic penalties can exponentially delay discovery in a deliberately
constructed finite sampling family. The sign of the effect depends on how
cost aligns with validity, and its magnitude depends on the penalty range.
The UGLY-SPACE Conjecture proposes that this mechanism matters for a
substantial family of open problems. We recommend occasionally checking
whether the hideous candidate is
correct before commissioning a nicer one.

\bibliographystyle{plain}
\bibliography{references}
\end{document}
